% =====================================================================
%  DROP-IN BLOCKS for spacetime_from_non-computability.tex
%  Drafted 2026-07-04 per JS's three promotions:
%   (1) relation-creation is experienced as acceleration -- promoted
%       from implication to its own paragraph;
%   (2) the Unruh effect as the framework's ledger-conversion, and as
%       Wigner's friend for accelerated observers;
%   (3) reversible continuity (sec:continuity) as the reason space is
%       two-way and time is not.
%  Plus (4) the section summary sentence updated to carry the new
%  yields. Four blocks, each with its own paste anchor, so nothing
%  collides with your current Overleaf state (including the
%  entropic-gravity paragraph if already pasted).
%
%  CITATIONS -- existing keys used: Penrose1996 (ambient, not re-cited),
%  DiBiagioRovelli2021, plus internal refs sec:comp-distance,
%  sec:continuity, sec:fanout-boundary, sec:wigners-friend, sec:jacobson.
%  ONE new key: unruh1976notes (RIS in spacetime_additions_refs_2026-07-04.ris);
%  swap to your canonical key if you already hold Unruh 1976.
%
%  PRECISION DECISIONS (encoded so the drafting history is clear):
%  - The "timelike converted to spacelike" claim is stated in the
%    framework's COMPUTATIONAL ledger (sec:comp-distance's own
%    definitions), with an explicit guard that the invariant metric
%    interval is untouched -- the observer-relative object is the
%    recordable/reversible partition, which is what the Unruh effect
%    exhibits. This honours the intent and is referee-proof.
%  - "Back and forth in space" is phrased as revisitability (a place can
%    be returned to; a moment is a stratum of records and cannot),
%    avoiding the naive reading that spatial motion is time-free.
%  - The equivalence-principle clause is a resonance, not a derivation;
%    all three additions are flagged structural, matching the section's
%    register.
%  - Hawking is deliberately implicit (via the Jacobson stitch);
%    Fulling/Davies omitted to keep the citation load minimal.
%  REVISED 2026-07-04 (same day, JS review):
%  - "interconvertible" -> "rotated into" (the tipping IS the rotation);
%  - the Chaitin mechanism added to Block 2: distance conversion =
%    addition/subtraction of information, completing sec:comp-distance's
%    own bits-added definition (reuses existing key Chaitin1998);
%  - BLOCK 5 added: the GR-internal derivation of d <= 3 via
%    full-circle cone-tipping (Penrose, Shadows of the Mind, p. 382 --
%    attribution corrected from The Emperor's New Mind per JS), with
%    the parallels-not-identity caveat stated in print and the
%    convergence point named: extreme tipping = closed timelike curve =
%    the object the computational definition counts. New key
%    Penrose1994 (keyed to match the existing Penrose1996 pattern; page
%    382 pinned by JS and carried as a biblatex postnote).
% =====================================================================

% ---------------------------------------------------------------------
%  BLOCK 1 + BLOCK 2 -- paste together at the END of
%  \subsection{Recovering the Lorentz group} (sec:lorentz), immediately
%  after the paragraph ending "...its identification with
%  measurement-as-tipping is the framework's structural reading." and
%  before \subsection{The relationalist commitment}.
% ---------------------------------------------------------------------

One consequence of the identification has been implied throughout this section and deserves its own
paragraph, because the next paragraph turns on it. If measurement creates a relation between
observer and observed, and if --- relationally --- a new relation alters the causal structure, then
the tipping of the cone is not merely \emph{associated} with the creation of the relation: it is
what the creation of the relation is, seen geometrically. And a tipping of one's cone against the
congruence of one's neighbours already has a name in physics: acceleration. The framework's
statement, made explicit here, is that \emph{creating a new relation between observer and observed
is experienced as acceleration}. The resonance with the equivalence principle is exactly what it
should be: the principle makes acceleration and gravitation locally indistinguishable, Penrose's
association --- run in this section's direction, reduction driving the tipping --- makes state
reduction and cone-tipping one operation, and the framework's identification closes the triangle,
measurement and acceleration meeting where the principle says they must. The claim is structural:
the framework's reading of the cone-tipping established above, not an independent derivation.

Taken literally, the identification has a measurable face, and it is a familiar one: the Unruh
effect, read in this section's own ledger. Recall the definitions of \S\ref{sec:comp-distance}:
timelike distance counts computable steps, spacelike distance counts non-computable bits. What an
observer can compute --- what their \textsc{fanout}s can convert into classical records --- is fixed
by their causal structure, and the previous paragraph has just made that structure
relation-dependent. An accelerated observer's cones are tipped, and with them the boundary of the
recordable: relations that for the inertial observer lie wholly in the reversible, pre-record sector
--- the entanglement structure of the vacuum, non-computable for them, spacelike in the
computational sense --- fall, for the accelerated observer, partly on the recordable side. The mechanism is the one
Chaitin identified for formal systems \autocite{Chaitin1998}: a proposition undecidable relative to
an axiom system becomes decidable when information is added to the axioms --- and
\S\ref{sec:comp-distance} defined spacelike distance as exactly the bits whose addition closes the
gap. Creating a relation is adding information to the observer's axiom base, so timelike and
spacelike distance are converted into one another by the addition and subtraction of information,
and acceleration is the continuous version of that addition. A
relation newly recordable is a detection. Particles appear from nowhere: an accelerated detector in
the inertial vacuum registers a thermal bath at $T = \hbar a / 2\pi c k_{\mathrm{B}}$
\autocite{unruh1976notes}. Two guards keep the reading honest. Nothing here reclassifies the
invariant interval --- timelike and spacelike separation in the emergent metric are Lorentz
invariants, and no event crosses that classification; what is observer-relative is the computational
partition, which relations a given observer's horizon renders recordable. In the ledger of
\S\ref{sec:comp-distance} that partition \emph{is} the timelike/spacelike distinction, so in the
framework's own units timelike distance is \emph{rotated} into spacelike, and spacelike into
timelike, by a change of observer --- the tipping of the cone is literally that rotation --- and the
detector's clicks are the rotation made audible. And the effect is established physics
while the reading is the framework's, flagged structural. What gives the reading its content is a
parallel: the inertial and the accelerated observer stand to the vacuum exactly as Wigner and his
friend stand to the laboratory (\S\ref{sec:wigners-friend}) --- one holds a coherent relation and no
event, the other holds a record, and the facts are relative on both sides in precisely Rovelli's
sense \autocite{DiBiagioRovelli2021}. Since the previous paragraph identifies acceleration with
relation-creation, the parallel is not an analogy but the same structure in gravitational dress, the
horizon standing where the laboratory wall stood: \emph{the Unruh effect is Wigner's friend for
accelerated observers}. And two subsections of this section that seemed independent here meet: the
temperature Jacobson's construction requires across every local Rindler horizon
(\S\ref{sec:jacobson}) is exactly this temperature --- the $T$ of the Clausius relation is the $T$
of the tipped cone, supplied by the very relation-creation with which this subsection began.

% ---------------------------------------------------------------------
%  BLOCK 3 -- paste at the END of \subsection{Timelike and spacelike
%  distance as computation} (sec:comp-distance), immediately after the
%  paragraph ending "...the identification is motivated by reproducing
%  what nature does, not proved." and before \subsection{Three
%  dimensions of space}.
% ---------------------------------------------------------------------

The split just defined explains the most familiar asymmetry in physics: that one can move back and
forth in space but not in time. \S\ref{sec:continuity} identifies the pre-measurement regime with
continuous reversibility --- its transformation group is connected, and every motion in it is
undoable. Spacelike separation, being the non-computable sector, is exactly the part of the
relational structure on which no record has been written; motion across it is motion within the
reversible regime, and it inherits that regime's continuity and its two-sidedness: a place can be
revisited because going and returning writes no record that forbids the return. Timelike separation
is the record chain itself --- proper time counts completed irreversible loops, as above --- and a
moment cannot be revisited because a moment \emph{is} a stratum of records, and returning would
require unwriting them, which is what \textsc{fanout} forbids (\S\ref{sec:fanout-boundary}). The
asymmetry between space and time is therefore not an extra postulate about the manifold: it is the
pre-/post-\textsc{fanout} asymmetry read as geometry --- space is what the reversible group feels
like from inside, time is what the record chain feels like --- and the arrow-of-time discussion
later in this section is the same asymmetry read thermodynamically. The identification is the
framework's structural reading of the split, offered in the same register as the rest of this
subsection.

% ---------------------------------------------------------------------
%  BLOCK 4 (optional but recommended) -- update the section summary so
%  the new yields are carried. REPLACE the sentence
%
%    "The spacetime arc therefore yields, from one definition of
%     spacelike distance as non-computability: the signature, three
%     spatial dimensions, the Lorentz group, the dynamical relational
%     metric, a reading of the holographic principle, and the presence
%     of non-locality within general relativity."
%
%  WITH:
% ---------------------------------------------------------------------

% The spacetime arc therefore yields, from one definition of spacelike
% distance as non-computability: the signature, three spatial dimensions,
% the Lorentz group, the dynamical relational metric, a reading of the
% holographic principle, the presence of non-locality within general
% relativity, the identification of relation-creation with acceleration
% and of the Unruh effect with Wigner's friend, and the space--time
% mobility asymmetry as the \textsc{fanout} asymmetry read as geometry.
%
%  ...and in the following sentence's conjectural list, after "the
%  general-relativistic non-locality", insert: ", the acceleration and
%  Unruh readings, and the mobility asymmetry" -- so the new items are
%  explicitly on the conjectural side of the ledger with the others.

% ---------------------------------------------------------------------
%  BLOCK 5 -- paste at the END of \subsection{Three dimensions of
%  space} (sec:three-dimensions), immediately after the paragraph
%  ending "...are named open problems." and before
%  \subsection{Recovering the Lorentz group}.
% ---------------------------------------------------------------------

A scoping remark, and then a second route to the same count. The computational model of spacetime
proposed here has many detailed parallels with the real thing, but a parallel is not an identity,
and the identity is not claimed in this paper: establishing it is work deferred to the second paper.
What can be claimed now is that the parallel is productive in a specific direction --- it suggests a
derivation of the dimension count carried out entirely within general relativity, with no
computational premise. Acceleration and gravitation tip the light cone, and the tipping can be
carried to the extreme: Penrose's discussion runs the cone through a full circle
\autocite[382]{Penrose1994}, at which point the tipped direction closes on itself --- which is to
say, the extreme of tipping is a closed timelike curve, the very object whose closure the
computational definition of \S\ref{sec:comp-distance} counts in bits. The two routes are one route
entered from opposite ends. Now demand of general relativity only this: that observers accelerating
with respect to one another, and with respect to unaccelerated observers, retain a consistent
conception of the distance between them. Each observer's tipping traces a path on a sphere of
directions, and mutual consistency requires the associated rotations to compose without regard to
bracketing --- to form a group. Rotations of $S^{1}$ and $S^{3}$ do, forming $U(1)$ and $SU(2)$ as
above; rotations of $S^{7}$ do not, being non-associative, so observers tipped across an
$S^{7}$-worth of directions could not agree on distance even pairwise: the same failure that forbade
$S^{7}$ a metric above forbids it here a shared ruler. The largest sphere admitting a consistent
conception of distance under mutual acceleration is therefore $S^{3}$, and space has at most three
dimensions; with the floor argued above --- the abelian $S^{1}$ alone cannot carry the required
structure --- the count is exactly three. The status matches the subsection's. The tipping and its
extreme are general relativity's own \autocite{Penrose1994}; the consistency demand is
elementary; the associativity filter is the algebraic fact used throughout; the reading of composed
tippings as motion on the parallelisable spheres is the structural step, shared with the
computational route. The convergence of the two routes on the same endpoint is offered as evidence
for the parallel, while the identity it gestures at remains the second paper's burden.
